
A systematic review of over 400 papers on AI alignment found that the vast majority focus exclusively on AI-to-human alignment: reward model quality, instruction following accuracy, and preference prediction \citep{Shen2024}. This unidirectional focus risks optimizing AI systems that are highly responsive to user preferences but may undermine users' ability to evaluate AI outputs critically.

The Bidirectional Alignment Framework \citep{Shen2024} proposes $Human\rightarrow AI$ alignment as a second necessary axis, encompassing human agency preservation, critical evaluation, and appropriate reliance on AI systems. However, this framework remains conceptual. No existing benchmark operationalizes bidirectional alignment with computable metrics, and no prior study has examined whether existing preference datasets contain extractable signals related to human agency preservation.

This work addresses a methodological question: can surface-level linguistic patterns serve as valid proxies for functional agency preservation? Four patterns are examined: clarifying questions, option enumeration, epistemic hedging, and explicit deferral. These patterns are motivated by HumanAgencyBench \citep{HumanAgencyBench2024}, which identifies six dimensions of agency-preserving AI behavior.

The Bidirectional Alignment Index (BAI) is introduced as a composite score computed from these four agency proxies applied to preference dataset responses. Experiments yield a combination of positive and negative results. BAI is highly extractable (mean AUROC 0.98 across proxies) and forms a statistically independent representational dimension (AUROC 0.99 after adversarial gradient reversal). However, semantic coherence analysis reveals that high-BAI responses cluster by generic conversational patterns rather than interpretable agency-preserving vocabulary.

The contributions are threefold. First, agency proxies are extracted from HH-RLHF and RewardBench preference data and their detection reliability is quantified. Second, adversarial probing is validated as a method for demonstrating representational independence between BAI and reward signals. Third, a negative result is established: surface proxies capture syntactic patterns rather than semantic agency content, indicating that future operationalizations require richer approaches.
